\documentclass[10pt]{article}
\usepackage[a4paper,margin=25mm]{geometry}
\usepackage{amsmath,amssymb}
\usepackage[T1]{fontenc}
\title{The stated numerical ratio in Conjecture 996 is false}
\author{Local AI-assisted solution draft}
\date{3 October 2026}
% Compact consistent title for a short mathematical note.
\makeatletter
\renewcommand{\maketitle}{\begin{center}
{\Large\bfseries\@title\par}\vspace{0.6em}
{\small\@author\par\@date}
\end{center}\vspace{0.5em}}
\makeatother
\setlength{\emergencystretch}{3em}
\usepackage{url}
\begin{document}
\maketitle
\paragraph{Source.}
Conjecture \texttt{00000000996} in The Last Math Competition was read at
the source commit
\begin{center}\small
\path{efab34b80a63963991d6c7ed625442a89a328a44}.
\end{center}
The statement explicitly specifies two
spectral dimensions $4/3$ and $7/4$ and claims their ratio is $7/9$.

\paragraph{Disproof.}
The asserted numerical identity fails in exact rational arithmetic:
\[
\frac{4/3}{7/4}=\frac{4}{3}\frac{4}{7}=\frac{16}{21}\ne\frac79,
\]
since $16\cdot9=144$ whereas $7\cdot21=147$.
Reading ``ratio between'' in the opposite order does not help:
\[
\frac{7/4}{4/3}=\frac{21}{16}\ne\frac79.
\]
Therefore the conjunction of the dimensions as written and the claimed
ratio cannot hold. This disproof does not require a claim about whether the
two individual spectral-dimension values are themselves correct.

\paragraph{Formal verification and scope.}
The self-contained Lean 4.19.0 file \texttt{Main.lean} uses the standard
library's exact rational type, \texttt{Std.Internal.Rat}. It defines the two
dimensions and the claimed ratio directly from the numerals in the source,
proves the forward and reverse ratio values, and proves
\path{Conjecture996.conjecture996_false}. The reversed reading is also
formally refuted. All reductions are kernel-checked, with no proof
placeholders, custom axioms, or native decision procedures. The only
reported foundational axiom is standard propositional extensionality.

This is a disproof of the statement's explicit arithmetic equality. It
establishes no additional percolation or spectral-dimension theorem. If the
organizers intended a different ratio or different values, that would be a
revised conjecture.

\paragraph{Status.}
This is a locally checked draft, not an organizer-accepted result. Prior
submission status must be refreshed before publication.
\end{document}
